\chapter*{Resumen} % si no queremos que añada la palabra "Capitulo"
\addcontentsline{toc}{chapter}{Resumen}% si queremos que aparezca en el índice
\markboth{RESUMEN}{RESUMEN} % encabezado
\lhead[\thepage]{\rightmark}
\rhead[\leftmark]{\thepage}
\pagenumbering{Roman} %comenzar la numeracion de paginas en numeros romanos
\setcounter{page}{3}


En este trabajo se aborda el problema fundamental de unificar la mecánica cuántica con la relatividad especial mediante el desarrollo riguroso de ecuaciones de onda relativistas. Se inicia con un análisis crítico de la \textbf{ecuación de Klein-Gordon}, destacando sus limitaciones: densidades de probabilidad negativas y soluciones de energía negativa, que la hacen físicamente inconsistente para describir partículas individuales.

Posteriormente, se deriva la \textbf{ecuación de Dirac} desde primeros principios, demostrando cómo su formulación matricial de primer orden resuelve los problemas de Klein-Gordon. El estudio enfatiza:

\begin{enumerate}
    \item El rol del \textbf{álgebra de Clifford} ($\{\gamma^\mu, \gamma^\nu\} = 2g^{\mu\nu}I$) en garantizar covarianza Lorentz.
    \item La emergencia natural del \textbf{espín} como grado de libertad intrínseco.
\end{enumerate}

Ademas mostramos que la forma de primer orden de la ecuación de Schrödinger puede obtenerse a partir de la ecuación de Dirac en el límite no relativista.Para este propósito, se introducen matrices que obedecen ciertas propiedades. Además, estudiamos el problema de una partícula con espín hacia arriba incidente sobre una barrera de potencial escalón y potencial de barrera finita y mostramos que se recuperan los resultados conocidos de la mecánica cuántica. Con base en estas conclusiones, proponemos que la ecuación de primer orden hayada es el límite no relativista de la ecuación de Dirac y describe de manera más apropiada a las partículas de espín $1/2$ en el régimen no relativista.
\vspace{0.5cm}
\noindent

